% !TEX root = ../main.tex

\chapter{总结与展望}
\chaptereng{Summary and Prospect}


\section{全文总结}

\subsection{主要结论}

截至中期报告节点，已完成的主要工作包括：

第一，火焰炉结构设计——完成了三种不同参数的炉子设计方案和相应的热力学计算、冷态仿真、火焰光谱仿真，验证了设计方案在3000 K高温工况下的可行性。

第二，TDLAS测量系统搭建——建立了完整的1.4$\mu$m波段TDLAS实验系统，包括激光器（标定至1397.824 nm）、光路、探测（TLC2201 TIA）、采集（10 kHz）等硬件模块。

第三，数据处理算法开发——开发了FSR刻度$\to$波数转换$\to$Voigt拟合$\to$温度反演的完整数据处理链路，通过了合成光谱验证（误差$<$2\%）。

第四，初步实验验证——完成了平面炉TDLAS扫描实验，成功分辨出双吸收峰并重建了二维温度场。

\subsection{创新点}

针对非预混平面绝热火焰炉炉普遍存在的高温极端工况下寿命短、燃烧不均匀不稳定等问题，优化当前火焰炉的内部结构，通过添加水冷结构和改变喷嘴阵列结构、数量和间距等关键参数实现更均匀的燃烧场分布、更优的排放性能和更高的燃烧温度。
采用更高分辨率的TDLAS和PLIF实验方法，对火焰进行标定，能够更详细的建立以非预混平面绝热火焰炉炉作为标定源的火焰的温度场浓度场，对微扩散火焰结构进行更深入的研究，为后续标定提供一个可溯源的标准化标定条件。


\section{展望}

下一步工作计划包括：进行平面炉氢氧火焰完整扫描实验并走通全流程温度反演；完成多工况系统实验、制备多通道探测器阵列、完成1分16光路实验和PLIF实验。